% tikz_capscheme.tex -- fig:cap-scheme, the cap argument of sec:cap-theorem.
%   (a) The plane x_2 = x_3 = 0 of the frame of four orthogonal roots, the
%       deviating root being e_0.  B = {|x_0|+|x_1|+|x_2|+|x_3| <= 2} meets it in
%       the square with vertices (+-2,0), (0,+-2); Q meets it in the pentagon
%       (2,0), (1,1), (-1,1), (-1,-1), (1,-1); the 24-cell in the square
%       [-1,1]^2.  The half-space x_0 >= 1 cuts the triangle (1,1), (2,0),
%       (1,-1) off both Q and B; in four dimensions it has volume 1/3
%       (lem:Q-structure).  A unit-distance half-space <x,w> >= 1, w at 25 deg,
%       meets the boundary of B at (0.32,1.68) and (1.3885,-0.6115).  To scale.
%   (b) The volume of the cap of B along two paths of unit normals from the
%       vertex direction e_0: to the midpoint of an edge (45 deg) and to the
%       centre of a facet (60 deg), from cap_paths.dat (cap_paths_data.py,
%       which evaluates lem:cap-formula and checks it by Monte Carlo).
% Test:   python3 tikz_overlap_check.py tikz_capscheme.tex
% Build:  pdflatex tikz_capscheme.tex && pdftoppm -r 500 -png -singlefile tikz_capscheme.pdf ../figures/fig_tikz_capscheme
\documentclass[border=4pt]{standalone}
\usepackage{tikz}
\usetikzlibrary{calc,arrows.meta,patterns}
\usepackage{pgfplots}
\pgfplotsset{compat=1.17}
\usepackage[T1]{fontenc}
\usepackage{lmodern}
\usepackage{amsmath,amssymb}
\DeclareMathOperator{\vol}{vol}
\definecolor{cblue}{HTML}{2A78D6}
\definecolor{corange}{HTML}{EB6834}
\definecolor{caqua}{HTML}{1BAF7A}
\definecolor{cink}{HTML}{0B0B0B}
\definecolor{cink2}{HTML}{52514E}
\definecolor{cink3}{HTML}{8B8A85}
\definecolor{cpanel}{HTML}{F4F3EF}
% tikzcheck.tex -- the hooks of tikz_overlap_check.py, input by the TikZ figures.
% Two ways to name what is checked.
%  - Explicit: \figboxes and \figlabels list named nodes; with \nolabels defined,
%    nodes in the style lbl are drawn invisibly.
%  - Automatic: \autocheck in the preamble gives every node of the picture,
%    pgfplots tick labels included, an alias chk<n>; with \nolabels defined, the
%    text of every node is drawn invisibly and its shape is kept.
% \checkfigure, at the end of the picture, writes the four corners of each
% checked node (so that rotated nodes are measured by their true extent) and
% those of the picture to the log when \checkbboxes is defined.
\tikzset{lbl/.style={}}
\ifdefined\nolabels\tikzset{lbl/.append style={opacity=0}}\fi
\newcount\chknodes
\newcommand\autocheck{% idempotent: a figure with \autocheck may also be run with --auto
  \ifdefined\chkauto\else
  \def\chkauto{}%
  \tikzset{every node/.append style={/utils/exec={\global\advance\chknodes by 1\relax},
                                      alias=chk\the\chknodes}}%
  \ifdefined\nolabels\tikzset{every node/.append style={text opacity=0}}\fi
  \fi}
\makeatletter
% a node counted but never kept (pgfplots typesets some nodes once to measure them)
% has no shape, and is skipped
\newcommand\chkifnode[2]{\pgfutil@ifundefined{pgf@sh@ns@#1}{}{#2}}
\makeatother
\newcommand\chkout[2]{% kind, node
  \path let \p1=(#2.south west), \p2=(#2.north east), \p3=(#2.north west), \p4=(#2.south east)
    in \pgfextra{\typeout{BBOX #1 #2\space\x1\space\y1\space\x2\space\y2\space\x3\space\y3\space\x4\space\y4}};}
\newcommand\checkfigure{%
  \ifdefined\checkbboxes
    \path let \p1=(current bounding box.south west), \p2=(current bounding box.north east)
      in \pgfextra{\typeout{PICBB \x1\space\y1\space\x2\space\y2}};
    \typeout{BORDER 4.0pt}%
    \ifdefined\chkauto
      \edef\chklast{\the\chknodes}%
      \foreach \nm in {1,...,\chklast} {\chkifnode{chk\nm}{\chkout{auto}{chk\nm}}}%
    \else
      \foreach \nm in \figboxes {\chkout{box}{\nm}}%
      \foreach \nm in \figlabels {\chkout{label}{\nm}}%
    \fi
  \fi}
\ifdefined\forceauto\autocheck\fi

\autocheck
\begin{document}
\begin{tikzpicture}[font=\small]
%% ------------------------------------------------------------------ (a)
\begin{scope}[x=1.45cm, y=1.45cm]
  % B minus Q
  \fill[cink3!18] (2,0) -- (0,2) -- (-2,0) -- (0,-2) -- cycle;
  % Q
  \fill[cblue!14] (2,0) -- (1,1) -- (-1,1) -- (-1,-1) -- (1,-1) -- cycle;
  % the 24-cell
  \fill[cblue!30] (-1,-1) rectangle (1,1);
  % the cap x_0 >= 1
  \fill[corange!55] (1,1) -- (2,0) -- (1,-1) -- cycle;
  % outlines
  \draw[cink2, thin, dashed] (2,0) -- (0,2) -- (-2,0) -- (0,-2) -- cycle;
  \draw[cblue!80!black, thick] (2,0) -- (1,1) -- (-1,1) -- (-1,-1) -- (1,-1) -- cycle;
  \draw[cink, thin] (1,-1) -- (1,1);
  % the unit circle
  \draw[cink2, densely dotted] (0,0) circle (1);
  % the tilted half-space <x,w> = 1, w at 25 degrees, tangent to the circle at w
  \draw[caqua!70!black, thick] ($(25:1)+(115:1.95)$) -- ($(25:1)+(-65:1.45)$);
  \draw[caqua!70!black, -{Stealth[length=5pt]}] (0,0) -- (25:1);
  \fill[cink] (25:1) circle (0.9pt);
  % axes
  \draw[-{Stealth[length=5pt]}, cink2] (-2.35,0) -- (2.45,0);
  \draw[-{Stealth[length=5pt]}, cink2] (0,-2.3) -- (0,2.35);
  \node[anchor=north, inner sep=2pt] at (2.42,-0.04) {$x_0$};
  \node[anchor=east, inner sep=2pt] at (-0.04,2.32) {$x_1$};
  % labels
  \node[inner sep=1.5pt] at (-1.42,0.3) {$B$};
  \node[inner sep=1.5pt] at (-0.55,-0.45) {$\mathcal V$};
  \node[inner sep=1.5pt] at (1.36,0.33) {$\tfrac13$};
  \node[anchor=north west, inner sep=1.5pt] at (1.04,-1.04) {$x_0=1$};
  \node[anchor=south west, inner sep=1.5pt, text=caqua!50!black] at (0.62,1.72) {$\langle x,w\rangle=1$};
  \node[anchor=north west, inner sep=1.5pt, text=caqua!50!black] at (0.62,0.36) {$w$};
\end{scope}
%% ------------------------------------------------------------------ (b)
\begin{axis}[name=A, at={(5.35cm,-2.75cm)}, anchor=south west, width=7.8cm, height=6.2cm,
    xmin=0, xmax=62, ymin=0, ymax=0.39,
    axis lines=left, tick style={cink2}, tick label style={font=\footnotesize},
    xtick={0,15,30,45,60}, ytick={0.1,0.2,0.3},
    yticklabels={$0.1$,$0.2$,$0.3$},
    xlabel={angle of $u$ from $e_0$ (degrees)}, ylabel={$\vol\bigl(B\cap\{\langle x,u\rangle\ge1\}\bigr)$},
    xlabel style={font=\footnotesize}, ylabel style={font=\footnotesize},
    clip=false]
  \addplot[cink3, dashed, domain=0:62, samples=2] {1/3};
  \addplot[corange, very thick] table[x=t, y=c1] {cap_paths.dat};
  \addplot[cblue, very thick] table[x=t, y=c2] {cap_paths.dat};
  \node[anchor=south west, inner sep=1.5pt, font=\footnotesize] at (axis cs:33,0.339) {$\tfrac13$, attained only at $u=e_0$};
  \node[anchor=south west, inner sep=1.5pt, font=\footnotesize, text=corange!80!black] at (axis cs:46,0.232) {to an edge};
  \node[anchor=west, inner sep=1.5pt, font=\footnotesize, text=cblue!80!black] at (axis cs:51,0.118) {to a facet};
\end{axis}
\node[anchor=north, inner sep=2pt] at ($(A.outer south)+(0,-0.02cm)$) {(b)};
\node[anchor=north, inner sep=2pt] at ($(0,0)!(A.outer south)!(0,1)$) {(a)};
\checkfigure
\end{tikzpicture}
\end{document}
